\section{Arbitrary thin smoothings on the actual local quotient}
\label{sec:local-variable-profiles}

The local theorems MA-LA-014 and MA-CQ-003--004 were first
proved with fixed profiles.  Their actual Fourier covariance
and trace-class compact-input and compact-output kernels
remove that restriction.  No bounds on derivatives of the
varying profiles will be imposed.

Use one of the actual rational or character local completions
\(\mathscr H_\beta\) constructed in MA-LA-011 or MA-CQ-002,
with the specified finite set \(S\), coherent additive character,
primitive finite factors, and probability compact Haar measure.
Let \(H_\beta:\mathscr H_\beta\to\mathscr H_{\bar\beta}\)
be its actual unitary Fourier map.  Thus
\(H_{\bar\beta}H_\beta=1\), including the finite Gauss phases
that cancel the odd real Fourier square.
For norm profiles \(\rho,b\), put
\[
 T(\rho,b;h)
    =M_{\rho(|x|_\infty)}A_h
                      H_{\bar\beta}M_{b(|\xi|_\infty)}H_\beta,
 \qquad g(y)=h(e^y).
 \tag{MA-LA.80}
\]
The rational trivial case is included.  These are the actual
quotient operators; the finite conductor and different
dilations are part of \(H_\beta\) throughout.

\paragraph{MA-LA-015 (profile comparison and the full local source limit).}
For a bounded position profile supported in a compact norm
annulus and a bounded frequency profile, (MA-LA.80) is trace
class.  When the position profile is smooth, it has the same
actual nuclear lift into the flat unit-shell section as
MA-LA-014 or MA-CQ-003, for every fixed profile.

The relevant exact covariance is
\[
 H_\beta A_h=A_{h^\vee}H_\beta,\qquad
 A_hH_{\bar\beta}=H_{\bar\beta}A_{h^\vee},\qquad
 h^\vee(t)=h(t^{-1}),\quad g^\vee(y)=g(-y).
 \tag{MA-LA.81}
\]
There are fixed nonnegative functions
\(\omega_g^{\rm out}(r),\omega_g^{\rm in}(r)\to0\)
as \(r\downarrow0\), depending only on the test in its
logarithmic local Hilbert realization, with the following
property.  If a bounded log multiplier \(a\), of norm at most
one, is supported in any interval of radius \(r\le1\), then
\[
 \|M_aC_g\|_1\le\omega_g^{\rm out}(r),\qquad
 \|C_gM_a\|_1\le\omega_g^{\rm in}(r).
 \tag{MA-LA.82}
\]
The bounds are independent of the centre of the interval.
For supports in \(k\) such windows, the corresponding bound
is \(k\) times the right side.
More quantitatively, for \(\operatorname{supp}g\subset[-m,m]\),
\[
 \omega_g^{\rm out}(r)+\omega_g^{\rm in}(r)
       \le C_m\|g\|_{C^4}\sqrt r,\qquad
 \|g\|_{C^4}=\max_{0\le j\le4}\|g^{(j)}\|_\infty ,
 \quad0<r\le1 .
 \tag{MA-LA.86a}
\]
The constant depends on the fixed compact auxiliary cutoff
and \(m\), not on a moving profile or its derivatives.

Here is a general moving-profile version of the local limit.
Let \(\ell_n<u_n\), \(v_n\in\mathbb R\), and
\(\epsilon_n\to0\).  Let \(0\le\rho_n,b_n\le1\)
have the following logarithmic support and plateau properties:
\[
 \begin{aligned}
 \rho_n(e^y)&=0&&
       (y\notin[\ell_n-\epsilon_n,u_n+\epsilon_n]),\\
 \rho_n(e^y)&=1&&
       (\ell_n+\epsilon_n\le y\le u_n-\epsilon_n),\\
 b_n(e^y)&=1&&(y\le v_n-\epsilon_n),\\
 b_n(e^y)&=0&&(y\ge v_n+\epsilon_n).
 \end{aligned}
 \tag{MA-LA.83}
\]
The position profiles are smooth; the usual smooth ball
frequency profiles may be used, and the trace comparison
also permits bounded frequency multipliers.
No monotonicity is required within a transition window.
Suppose
\[
                u_n+v_n\longrightarrow+\infty,\qquad
                \ell_n+v_n\longrightarrow-\infty .
 \tag{MA-LA.84}
\]
For every compact smooth \(h\) with \(h(1)=0\), the actual
trace satisfies
\[
       \operatorname{Tr}T(\rho_n,b_n;h)
            \longrightarrow A_\infty(h)+B_{\beta,S}(h).
 \tag{MA-LA.85}
\]
Here the archimedean term, primitive directional prime
coefficients and conductor convention are exactly
MA-CQ.21--.23, or MA-LA.69--.70 in the rational sector.
If instead \(\ell_n+v_n\to+\infty\) and
\(u_n-\ell_n\to+\infty\), the limit is zero.

\begin{proof}
\emph{Actual finite-profile trace class and the small-window bounds.}
Use the exact logarithmic local unitary MA-CQ.19, or its
rational version.  It identifies \(A_h\) with \(C_g\) on
ordinary \(L^2(\mathbb R,dy)\), and norm multipliers with
their functions of \(e^y\).
For a smooth compact position multiplier,
the kernel of \(M_\rho C_g\) is smooth and supported in a
compact rectangle.  A double Fourier-series expansion on a
containing rectangle has rapidly summable coefficients.
It therefore gives an actual nuclear map into the smooth
compact output interval, with every derivative seminorm.
The finite unit-shell section puts that output into the
same flat Schwartz source as in MA-CQ.24.
Composition with the remaining bounded Fourier-cutoff
operator preserves this nuclear lift and its Hilbert trace.
No uniform bound for these LF derivative seminorms is
claimed when the profile varies.
For a merely bounded compact position multiplier, factor it
through a fixed smooth compact multiplier equal to one on
its support.  This still proves Hilbert trace class, but
does not assert a lift into a smooth section for a
discontinuous output profile.

Fix \(\chi\in C_c^\infty(\mathbb R)\) equal to one on
\([-2,2]\).  Both
\[
 \begin{aligned}
 K_g^{\rm out}&=M_\chi C_g\in\mathcal S_1,\\
 K_g^{\rm in}&=C_gM_\chi\in\mathcal S_1,\qquad
 P_r=M_{\mathbf1_{[-r,r]}},\\
 \omega_g^{\rm out}(r)&=\|P_rK_g^{\rm out}\|_1,\qquad
 \omega_g^{\rm in}(r)=\|K_g^{\rm in}P_r\|_1 .
 \end{aligned}
 \tag{MA-LA.86}
\]
Trace class follows from the same compact-rectangle
calculation.  The projections \(P_r\) converge strongly
to zero as \(r\downarrow0\).
Finite-rank approximation of a trace-class operator
therefore shows that both displayed functions tend to zero.
For a multiplier supported in \([-r,r]\), factor through
\(P_r\) and \(\chi\) to obtain (MA-LA.82).
Unitary log translation moves any centre to zero and
commutes with \(C_g\), proving that the bound is independent
of the centre.  Splitting a multiplier supported on finitely
many windows into bounded pieces proves the last assertion.

For (MA-LA.86a), choose containing rectangles for the two
smooth kernels with side lengths depending only on \(m\)
and the fixed \(\chi\).
Their normalized trigonometric bases \(e_j,f_k\) have
supremum norms bounded independently of \(j,k\).
Two integrations by parts in each variable give Fourier
coefficients
\[
 |a_{jk}|
     \le C_m\|g\|_{C^4}(1+j^2)^{-1}(1+k^2)^{-1}.
\]
There are no boundary terms because the kernels vanish
near the edges of the rectangles.
Only derivatives of \(g\) of total order at most four
occur.  Hence the coefficient matrix is absolutely
summable with sum at most \(C_m\|g\|_{C^4}\).
The rank-one Fourier expansion for \(K_g^{\rm out}\)
now gives
\(\|P_rK_g^{\rm out}\|_1
 \le\sum_{j,k}|a_{jk}|\|P_re_j\|_2\|f_k\|_2
 \le C_m\|g\|_{C^4}\sqrt r\).
For \(K_g^{\rm in}P_r\), restrict the input basis instead.
This proves the stated rate for both moduli.

\emph{Fourier covariance with the correct dual sector.}
For normalized norm dilation
\(\mathcal D_tu(x)=t^{-1/2}u(t^{-1/d_\infty}x)\),
additive Fourier scaling gives
\(\mathcal F_\infty\mathcal D_t
 =\mathcal D_{t^{-1}}\mathcal F_\infty\).
The exponent \(1/2\) is the unitary norm normalization;
the real or complex additive measure scales by \(t\).
Every actual finite Fourier factor in MA-CQ.10--.11
is a function of commuting norm dilations, including
its conductor/different shift and Gauss scalar.
It consequently commutes with every \(A_h\).
Integrating the dilation identity against \(h(t)\,dt/t\)
proves (MA-LA.81).  The second identity can also be
obtained from the first using \(H_{\bar\beta}H_\beta=1\).
This proof keeps the intermediate \(\bar\beta\) sector;
it does not replace the odd Fourier map by an involution
before its finite phases have been included.

For two frequency profiles the actual operator difference
thus factors as
\[
 \begin{split}
 T(\rho,b_1;h)-T(\rho,b_2;h)
   &=M_\rho H_{\bar\beta}
           A_{h^\vee}M_{b_1-b_2}H_\beta,\\
 \|T(\rho,b_1;h)-T(\rho,b_2;h)\|_1
   &\le\|C_{g^\vee}M_{(b_1-b_2)(e^\cdot)}\|_1 .
 \end{split}
 \tag{MA-LA.87}
\]
For two position profiles the corresponding bound is
\[
 \|T(\rho_1,b;h)-T(\rho_2,b;h)\|_1
       \le\|M_{(\rho_1-\rho_2)(e^\cdot)}C_g\|_1 .
 \tag{MA-LA.88}
\]
The multipliers and Fourier maps outside the indicated
trace-class factors have norm at most one.
These are operator trace-norm bounds, not a formal
cyclic rearrangement of an unconstructed trace.
In particular no derivatives of either varying profile
occur, and no convolution has been commuted past a
Fourier cutoff.

\emph{Comparison with fixed smooth profiles.}
Fix \(0<\eta<\min(1/4,\log2)\).
Choose smooth nonincreasing \(A_\eta,B_\eta\) on
\([0,\infty)\), between zero and one, equal to one on
\([0,1]\) and zero on \([e^\eta,\infty)\).
Define
\[
 \begin{aligned}
 \rho_n^\eta(t)
   &=A_\eta(t/e^{u_n})-A_\eta(t/e^{\ell_n}),\\
 b_n^\eta(t)&=B_\eta(t/e^{v_n}),\qquad
 T_n^\eta=T(\rho_n^\eta,b_n^\eta;h).
 \end{aligned}
 \tag{MA-LA.89}
\]
These are the fixed-profile operators of MA-LA-014 or
MA-CQ-003, with upper position radius \(e^{u_n}\),
lower radius \(e^{\ell_n}\), and frequency radius
\(e^{v_n}\).  The necessary separation of the two
position endpoints holds eventually, because
(MA-LA.84) implies \(u_n-\ell_n\to\infty\).

The difference \(\rho_n-\rho_n^\eta\) is supported
in two log windows about \(\ell_n,u_n\), each of
radius at most \(\epsilon_n+\eta\).
The difference \(b_n-b_n^\eta\) is supported in one
log window about \(v_n\) with the same radius.
All these differences have absolute value at most one.
Equations (MA-LA.82), (MA-LA.87)--(MA-LA.88) give
\[
 \|T(\rho_n,b_n;h)-T_n^\eta\|_1
   \le2\omega_g^{\rm out}(\epsilon_n+\eta)
               +\omega_{g^\vee}^{\rm in}(\epsilon_n+\eta).
 \tag{MA-LA.90}
\]
For \(\operatorname{supp}g\subset[-m,m]\), (MA-LA.86a)
bounds this error by
\(C_m\|g\|_{C^4}\sqrt{\epsilon_n+\eta}\).
This is a quantitative comparison between actual profile
realizations.  It does not assert a rate for the remaining
fixed-profile geometric limit.
For each fixed \(\eta\), the already proved fixed-profile
theorem gives \(\operatorname{Tr}T_n^\eta\to
A_\infty+B_{\beta,S}\), by (MA-LA.84).
The limit is independent of \(\eta\).
Taking the upper limit in \(n\) in (MA-LA.90), then
\(\eta\downarrow0\), proves (MA-LA.85).
The fixed-profile zero limit under the second stated
regime gives the other conclusion in exactly the same
way.  Only differences of profile realizations are
controlled in trace norm here; the varying operators
\(T(\rho_n,b_n;h)\) themselves are not asserted to
converge in \(\mathcal S_1\).
\end{proof}

\paragraph{The printed source smoothing conditions and the full global comparison.}
For the position cutoff in Masullo, Definition~9.3.3.1,
printed p.~427, take a smooth \(0\le\rho_n\le1\) supported
in \([1/X_n,X_n]\), equal to one on
\([1/X_n+\delta_n,X_n-\delta_n]\), with its stipulated
multiplicative symmetry if desired.
For the actual PW frequency ball, take \(0\le b_n\le1\)
equal to one up to \(Y_n-\delta_n\), and zero from
\(Y_n\) onward.  An outward thin smoothing is covered
as well.  Inward smoothing satisfies (MA-LA.83) with
\[
 \begin{gathered}
 \ell_n=-\log X_n,\quad u_n=\log X_n,\quad v_n=\log Y_n,\\
 \epsilon_n=
 \max\{\log(1+\delta_nX_n),
             -\log(1-\delta_n/X_n),
             -\log(1-\delta_n/Y_n)\}\longrightarrow0 .
 \end{gathered}
 \tag{MA-LA.91}
\]
For positive \(\delta_n\) all denominators here are valid
eventually.  The printed conditions
\(\delta_nX_n\to0\), \(\delta_nY_n\to0\) on p.~519,
together with \(X_n,Y_n\to\infty\), imply this conclusion.
Independent position and frequency transition widths
are equally allowed when their corresponding logarithmic
widths tend to zero.
The bounds permit arbitrarily rapid smooth oscillations
inside each transition interval.

The support convention at frequency is not inferred
from the annular definition on p.~427.
Masullo's Step~II on printed p.~516 explicitly specifies
the Fourier subset \(\{|x|_{F,S}\le\lambda'\}\);
the displayed composition on pp.~516--517 is its
Fourier-conjugated ball multiplier.
The footnote on p.~514 separately discusses the
one-sided position intention and an independently
moving lower endpoint.  The present theorem extends
the actual repaired annular position operator and
that specified frequency ball, retaining a smooth
flat section at each finite annular cutoff.

If also \(Y_n/X_n\to0\), then (MA-LA.84) holds and
the local limit is (MA-LA.85), for every such choice
of smooth source profiles.  The global source-profile
limit has already been proved in MA-FT-001 at every
polynomial weight and for every compact smooth
\(g\) with \(g(0)=0\).
Thus its paired local correction (MA-FT.7), its
full-strip comparison (MA-FT.17), and its exact
arithmetic gate MA-FT-002 remain valid with these
arbitrary thin source local smoothings, in place
of the initially fixed local profiles.
For a character the level remains \(Y=Q\Lambda\),
and \(S\) includes the conductor/different places
and all prime-power contributions to the test.
No finite Euler factor, Gauss phase or conductor
dilation is removed in making this extension.

\paragraph{MA-LA-016 (an annular frequency cutoff is a different operator).}
Keep the same source position cutoffs, but now replace
the PW frequency ball by a smooth annular frequency
profile \(\widetilde b_n\), bounded between zero and one,
supported in \([1/Y_n,Y_n]\), and equal to one on
\([1/Y_n+\delta_n,Y_n-\delta_n]\).
Assume
\[
 X_n,Y_n\to\infty,\qquad Y_n/X_n\to0,\qquad
 \delta_nX_n\to0,\qquad\delta_nY_n\to0 .
 \tag{MA-LA.92}
\]
Then its actual local operator is trace class, with
a nuclear flat-section lift for every smooth
position profile, and
\[
             \operatorname{Tr}T(\rho_n,\widetilde b_n;h)
                         \longrightarrow0
                         \qquad(h(1)=0).
 \tag{MA-LA.93}
\]
This is not the frequency operator specified on
Masullo p.~516.

\begin{proof}
For fixed \(0<\eta<\min(1/4,\log2)\), use the same
position reference as before and put
\[
 \widetilde b_n^\eta(t)
        =B_\eta(t/Y_n)-B_\eta(Y_nt).
 \tag{MA-LA.94}
\]
For large \(n\) this is bounded between zero and one.
Its difference from \(\widetilde b_n\) is supported
in two shrinking log windows about
\(-\log Y_n,\log Y_n\), up to the fixed reference
width \(\eta\).
The lower relative source width is now
\(\log(1+\delta_nY_n)\to0\).
The same proof of (MA-LA.90), with two frequency
windows instead of one, makes the trace error
arbitrarily small as first \(n\to\infty\), then
\(\eta\to0\).

Linearity of the actual bounded frequency multiplier
gives the exact fixed-reference decomposition
\[
 \begin{split}
 T(\rho_n^\eta,\widetilde b_n^\eta;h)
    ={}&T(\rho_n^\eta,B_\eta(\,\cdot/Y_n);h)\\
       &-T(\rho_n^\eta,B_\eta(Y_n\,\cdot);h).
 \end{split}
 \tag{MA-LA.95}
\]
Both terms are the actual ball operators already
constructed, with frequency radii \(Y_n\) and
\(1/Y_n\), respectively.
For the first, the two coupled products are
\(X_nY_n\to\infty\) and \(Y_n/X_n\to0\).
For the second, they are
\(X_n/Y_n\to\infty\) and \(1/(X_nY_n)\to0\).
The fixed-profile local theorem therefore gives
the same limit \(A_\infty(h)+B_{\beta,S}(h)\)
for both terms.
Their difference tends to zero, and the profile
comparison proves (MA-LA.93).
\end{proof}

This alternative changes an actually detectable trace.
For example, choose a nonzero nonnegative
\[
                   h\in C_c^\infty((1,\sqrt2)).
 \tag{MA-LA.96}
\]
Then \(h(1)=0\), every prime-power and reciprocal
prime-power contribution is zero, and the real
parity or complex archimedean integral in MA-CQ.21
is strictly positive.
The actual ball operator therefore has positive
limit \(A_\infty(h)\), whereas (MA-LA.93) gives
zero for the annular frequency alternative.
This test uses no zero-location information.
It is a scope counterexample to changing the
frequency support convention, not an attribution
of that changed convention to the printed PW map.

The trace framework is the one already justified
by the actual nuclear section and common Hilbert
trace in MA-LA-014 and MA-CQ-003.
The new ingredient is the derivative-free
trace-norm comparison (MA-LA.87)--(MA-LA.90).
It enlarges the proved cutoff class while leaving
the finite-length arithmetic deficiency of
MA-FT-002 explicit; it provides no further
zero-location or arithmetic positivity theorem.
